\subsection{流形的乘积}

5.6.1. 设 X 和 X' 是两个集合，且 \(c = (\mathrm{U}, \varphi, \mathrm{E})\)（分别地，\(c' = (\mathrm{U}', \varphi, \mathrm{E}')\)）是 X（分别地，X'）的一个图。三元组

\[
(\mathbf {U} \times \mathbf {U} ^ {\prime}, \varphi \times \varphi^ {\prime}, \mathbf {E} \times \mathbf {E} ^ {\prime})
\]

是乘积集 \(X \times X'\) 的一个图，记为 \(c \times c'\)。

5.6.2. 设 X 和 X' 是两个 C＾\(C^{r}\)\(C^{r}\)r 类流形。在乘积集 X × X' 上存在唯一的 C＾r 类流形结构，使得对于 X 的每个图 c 和 X' 的每个图 c'，c × c' 都是 X × X' 的一个图。赋予此结构后，称 X × X' 为流形 X 和 X' 的乘积流形。任意有限个流形的乘积类似地定义。流形 X × X' 的拓扑是 X 和 X' 的拓扑的乘积拓扑。对于 a ∈ X 和 b ∈ X'，有：

\[
\dim_ {(a, b)} (\mathrm{X} \times \mathrm{X} ^ {\prime}) = \dim_ {a} \mathrm{X} + \dim_ {b} \mathrm{X} ^ {\prime}.
\]

5.6.3. 设 X 和 X' 是两个流形，X × X' 是它们的乘积。令 a ∈ X，a' ∈ X'。典范投影

\[
\operatorname{pr} _ {1}: \mathrm{X} \times \mathrm{X} ^ {\prime} \rightarrow \mathrm{X} \quad \text { and } \quad \operatorname{pr} _ {2}: \mathrm{X} \times \mathrm{X} ^ {\prime} \rightarrow \mathrm{X} ^ {\prime}
\]

是态射。设 \(\pi_{i}\)（\(i=1,2\)）是它们在点 \((a,a')\) 处的切线性映射。映射

\[
(\pi_ {1}, \pi_ {2}): \mathrm{T} _ {(a, a ^ {\prime})} (\mathrm{X} \times \mathrm{X} ^ {\prime}) \rightarrow \mathrm{T} _ {a} (\mathrm{X}) \times \mathrm{T} _ {a ^ {\prime}} (\mathrm{X} ^ {\prime})
\]

是一个同构，从而可以把 X × X' 在 (a, a') 处的切空间与乘积 \(T_{a}(X) \times T_{a'}(X')\) 等同起来。

由此等同得到的从 \(\mathbf{T}_{a}(\mathbf{X})\) 到 \(\mathbf{T}_{(a,a^{\prime})}(\mathbf{X}\times \mathbf{X}^{\prime})\) 的嵌入，是从 \(a\) 到 \(x\mapsto (x,a^{\prime})\) 的态射 \(\mathbf{X}\) 在 \(\mathbf{X}\times \mathbf{X}^{\prime}\) 处的切线性映射；对于从 \(\mathbf{T}_{a^{\prime}}(\mathbf{X}^{\prime})\) 到 \(\mathbf{T}_{(a,a^{\prime})}(\mathbf{X}\times \mathbf{X}^{\prime})\) 的嵌入，也有类似结果。

5.6.4. 设 W、X 和 X' 是三个流形，且 \(f: \mathbf{W} \to \mathbf{X}\)、\(f': \mathbf{W} \to \mathbf{X'}\) 是两个映射。对于映射

\[
(f, f ^ {\prime}): \mathbf {W} \rightarrow \mathbf {X} \times \mathbf {X} ^ {\prime}
\]

要成为态射，其充分必要条件是 \(f\) 和 \(f'\) 都是态射（这说明了使用“乘积”一词的合理性，参见~集合论，第~IV~章，§ 2，第 4 号）。在此情形下，若 \(a\) 是 W 的一点，则有：

\[
\mathrm{T} _ {a} (f, f ^ {\prime}) = (\mathrm{T} _ {a} (f), \mathrm{T} _ {a} (f ^ {\prime})),
\]

其中考虑了上面所作的等同。

5.6.5. 设 \(f: \mathbf{X} \to \mathbf{Y}\) 和 \(f': \mathbf{X}' \to \mathbf{Y}'\) 是流形的态射。则 \(f \times f': \mathbf{X} \times \mathbf{X}' \to \mathbf{Y} \times \mathbf{Y}'\) 是一个态射。若 \(a \in \mathbf{X}\) 且 \(a' \in \mathbf{X}'\)，则有：

\[
\begin{array}{l} \mathrm{T} _ {(a, a ^ {\prime})} (f \times f ^ {\prime}) = \mathrm{T} _ {a} (f) \times \mathrm{T} _ {a ^ {\prime}} (f ^ {\prime}) \\ \mathrm{rg} _ {(a, a ^ {\prime})} (f \times f ^ {\prime}) = \mathrm{rg} _ {a} f + \mathrm{rg} _ {a ^ {\prime}} f ^ {\prime}. \end{array}
\]

5.6.6. 设 \(X_{1}\)、\(X_{2}\) 和 Z 是三个流形，f 是从 \(X_{1} \times X_{2}\) 到 Z 的态射。令 \(a \in X_{1}\)，\(b \in X_{2}\)。将 \(x \mapsto f(x, b)\)（分别为 \(y \mapsto f(a, y)\)）到 Z 的偏映射 \(X_{1}\)（分别为 \(X_{2}\)）的切线性映射记为 \(T_{(a,b)}^{1}(f)\)（分别为 \(T_{(a,b)}^{2}(f)\)）。若把 \(T_{(a,b)}(X_{1} \times X_{2})\) 与 \(T_{a}(X_{1}) \times T_{b}(X_{2})\) 等同，则有：

\[
\mathrm{T} _ {(a, b)} (f). (u, v) = \mathrm{T} _ {(a, b)} ^ {1} (f). u + \mathrm{T} _ {(a, b)} ^ {2} (f). v
\]

对于每个 \(u \in \mathrm{T}_a(\mathrm{X}_1)\) 和每个 \(v \in \mathrm{T}_b(\mathrm{X}_2)\)v  T＿b(X＿2)。

5.6.7.（“隐函数定理”）。在前一号的假设和记号下，再设 \(\mathrm{T}_{(a,b)}^{2}(f)\) 是双射。则存在 \(X_{1}\) 中 a 的一个开邻域 U 和 \(X_{2}\) 中 b 的一个开邻域 V，具有如下性质：对于每个 \(x \in U\)，存在 V 中唯一的一点 \(g(x)\)，使得 \(f(x, g(x)) = f(a, b)\)，且映射 g 是从 U 到 V 的态射。对于每个 \(x \in U\)，有

\[
\mathrm{T} _ {x} (g) = - \left(\mathrm{T} _ {(x, g (x))} ^ {2} (f)\right) ^ {- 1} \circ \mathrm{T} _ {(x, g (x))} ^ {1} (f).
\]

